\begin{resumo}

O ensino prático da Engenharia de Requisitos (ER) enfrenta desafios logísticos significativos para simular as ambiguidades e a complexidade sociotécnica inerentes ao processo real de elicitação. Embora a prova de conceito do agente conversacional ElicitBot tenha demonstrado potencial pedagógico inicial em sessões individuais, limitações de acoplamento arquitetural, volatilidade de estado em memória e crescimento cumulativo no consumo de \textit{tokens} impediam sua operação simultânea em turmas completas. Este trabalho teve como objetivo viabilizar técnica e funcionalmente a ferramenta ElicitBot, permitindo seu uso controlado como recurso pedagógico de apoio à prática de elicitação de requisitos em turmas da disciplina de Requisitos de Software. Para isso, adotou-se uma metodologia orientada pelos princípios da \textit{Design Science Research} (DSR), estruturada em quatro fases: diagnóstico e planejamento, refatoração e desenvolvimento, implantação e validação em campo, e avaliação e análise de resultados. A partir do diagnóstico do sistema herdado, da elaboração de oito Registros de Decisão Arquitetural (ADRs) e da modelagem financeira baseada em \textit{Cloud Economics}, a aplicação foi reestruturada em uma arquitetura cliente-servidor desacoplada (React, FastAPI e PostgreSQL) sob o padrão transacional \textit{Read-Release-Write} e implantada em servidor institucional conteinerizado. O motor conversacional evoluiu para uma Máquina de Estados Finita de 12 etapas articulada ao gerenciamento bipartido de memória com \textit{Context Caching}, ao controle de ciclo de vida de oito categorias de artefatos rastreáveis com suporte à intervenção humana (\textit{Human-in-the-Loop}) e à dupla validação de Requisitos Funcionais, registrando 0\% de falhas para até 50 usuários simultâneos nos testes de carga e resistência com a ferramenta \textit{Locust}. A validação empírica em sala de aula contou com a participação simultânea de 78 estudantes distribuídos em 14 equipes (processando 519 turnos conversacionais), a coleta de $N = 65$ avaliações individuais válidas ($\alpha = 0,877$) e a aplicação imediata de refinamentos corretivos na interface e no \textit{prompt}. Os resultados demonstraram mediana e moda iguais à nota máxima (5,0) em todos os 20 critérios avaliados (média global de $\mu = 4,49$), destacando-se 100,0\% de aprovação em adesão ao fluxo lógico, 96,6\% em rastreabilidade de requisitos, 90,8\% de satisfação com a interação e 83,1\% de potencial de adoção, acompanhados de taxa de reaproveitamento de \textit{cache} em 74,0\% dos turnos, latência mediana de 5,14 s e custo total faturado de R\$ 20,84 para atender as duas turmas avaliadas. Conclui-se que a articulação entre infraestrutura assíncrona persistente, engenharia de contexto orientada a estados e mecanismos de governança e controle pelo usuário viabiliza a adoção escalável e economicamente sustentável de agentes baseados em Modelos de Linguagem de Grande Porte (LLMs) no apoio ao ensino prático de Engenharia de Requisitos.

\vspace{\onelineskip}
    
\noindent
\textbf{Palavras-chave}: Engenharia de Requisitos. Inteligência Artificial. Modelos de Linguagem de Grande Porte. Educação Inteligente. ElicitBot.
\end{resumo}
